% The full proof is dimension84_obstruction.md; exact replay is adjacent.
\begin{theorem}[Dimension-84 adjoint rigidity]\label{thm:m3-dimension84}
No complex fusion category has dimension type $(1,4;2,2;6,2)$ and
$\mathcal C_{\rm ad}=\mathcal C$. Consequently the five tensors
$R8\_3\_76$, $R8\_3\_735$, $R8\_3\_738$, $R8\_3\_1422$, and
$R8\_3\_1805$ are not categorifiable.
\end{theorem}
\begin{proof}
By \cite[Theorem~1.2]{DongWeak}, a category of dimension $84$ is
solvable or group-theoretical. Every group of order $84$ is solvable:
its Sylow $7$ subgroup is normal and its quotient has order $12$.
Solvability is Morita invariant, so either alternative makes
$\mathcal C$ solvable. Its last cyclic-prime construction step
\cite[Definition~1.2, Proposition~4.4]{ENOSolvable} cannot be an
extension because $\mathcal C_{\rm ad}=\mathcal C$. It is therefore
an equivariantization $\mathcal C=\mathcal D^{C_2}$: the prime divides
the number four of invertibles. The base has dimension $42$ and is
group-theoretical by \cite[Theorem~9.2]{ENOSolvable}. It is integral,
since the orbit dimension formula makes its simple dimensions rational
algebraic integers. Exactly two base invertibles are fixed by $C_2$.
Let $g$ be the sign object of the equivariantization. If $n_2,n_6$
count its fixed simples of dimensions two and six, respectively, then
$n_2,n_6\in\{0,2\}$ and the base type is
\[
 (1,2+2n_2;\ 2,(2-n_2)/2;\ 3,2n_6;\ 6,(2-n_6)/2).
\]
Thus the four possible multisets of base dimensions are
$\{1,1,2,6\}$, $\{1^6,6\}$, $\{1,1,2,3^4\}$, and $\{1^6,3^4\}$.
For completeness, all groups of order $42$ are $C_7\rtimes C_6$
with action image of order $1,2,3,6$, and $C_7\rtimes S_3$ with
action image of order $1,2$. Exhausting their subgroups and double
cosets gives the following complete list of group-theoretical bases
of the four types. The first and third types never occur. Write
$F_{21}=C_7\rtimes C_3$ and $F_{42}=C_7\rtimes C_6$ for the
faithful semidirect products. Type $\{1^6,6\}$ occurs only for
$(G,H)=(F_{42},C_6)$ or $(F_{42},F_{42})$. Type $\{1^6,3^4\}$
occurs only for $G=F_{21}\times C_2$, with $|H|=3,6,21,42$,
or $G=F_{42}$, with $|H|=3,21$, up to conjugacy.
This finite audit includes arbitrary associators: every double-coset
stabilizer has squarefree order and trivial Schur multiplier, by
restriction/corestriction to its cyclic Sylow subgroups. Its projective
character degrees are therefore ordinary character degrees. The exact
subgroup and double-coset enumeration is part of the certificate.
In the second type the category is $\operatorname{Rep}(F_{42})$.
For $H=C_6$, the restriction $H^3(F_{42},\mathbb C^\times)\to
H^3(C_6,\mathbb C^\times)$ is an isomorphism: the remaining summand
$H^3(C_7,\mathbb C^\times)^{C_6}$ vanishes because the faithful
multiplier acts by its square. Thus the required trivial restriction
kills the associator, and the split exact-factorization construction
is $\operatorname{Rep}(F_{42})$; compare
\cite[Corollary~7.4]{EGODimPQ}.
In the fourth type the canonical adjoint subcategory is
$\operatorname{Rep}(F_{21})$. For $H=C_3$ or $F_{21}$, support in
the normal index-two subgroup $F_{21}$ gives a faithful $C_2$ grading
whose neutral category has type $(1,3;3,2)$. It is
$\operatorname{Rep}(F_{21})$ by \cite[Theorem~6.3]{EGODimPQ}.
Since $F_{21}$ is centerless, this neutral category equals its own
adjoint, forcing equality with the adjoint of the whole category.
For $G=F_{21}\times C_2$ and $H=C_6$ or $G$, coprime cohomology
decomposition instead gives
$\operatorname{Rep}(F_{21})\boxtimes\operatorname{Rep}(C_2)$,
with the same adjoint. These reductions do not assume a trivial
ambient associator.
A representation category of a squarefree-order group has a unique
fiber functor: the subgroup supporting a nondegenerate projective
representation must have square order
\cite[Section~3.4, Theorem~2 and Example~(i)]{OstrikModules}.
Hence every tensor autoequivalence of $\operatorname{Rep}(F_{21})$
or $\operatorname{Rep}(F_{42})$ comes from a group automorphism.
In the presentation $a^7=b^m=1$, $bab^{-1}=a^\lambda$, with
$m=3$ or $6$ and $\lambda$ of order $m$, an automorphism
$a\mapsto a^u$, $b\mapsto a^v b^w$ forces $\lambda^w=\lambda$.
It therefore fixes the abelianization and every linear character.
The second base type consequently has six fixed invertibles; in the
fourth type the canonical adjoint supplies at least three fixed
invertibles. Both contradict the required number two.
For each displayed source tensor, the certificate verifies all
dimension equations, full adjoint closure, and every possible sign
permutation directly in the given basis.
\end{proof}
\begin{remark}[A restricted dimension-660 candidate]\label{rem:m3-r731}
The residual ring $R8\_3\_731$ is integral of dimension $660$ and
has no nontrivial proper closed based subring. Its exact primary-three
witness excludes unitary categorification. By
\cite[Proposition~9.11]{ENOSolvable}, a hypothetical weakly
group-theoretical categorification would be a representation category
of a nonabelian finite simple group, hence unitary, a contradiction.
Thus any categorification of this ring would be nonunitary and not
weakly group-theoretical. Independently, the ordinary representation
ring scan finds no matching tensor among the $40$ groups of order
$660$. General complex categorification remains unresolved.
\end{remark}
